\documentclass[11pt,a4paper]{article}
\usepackage[T1]{fontenc}
\usepackage[utf8]{inputenc}
\usepackage{lmodern}
\usepackage[french]{babel}
\usepackage[a4paper,hmargin=2.2cm,top=2.1cm,bottom=2cm,headsep=0.6cm]{geometry}
\usepackage[dvipsnames]{xcolor}
\usepackage{fancyvrb}
\usepackage{enumitem}
\usepackage{titlesec}
\usepackage{fancyhdr}
\usepackage[most]{tcolorbox}
\usepackage{tikz}
\usepackage[hidelinks]{hyperref}

% --- Couleurs -----------------------------------------------------------
\definecolor{bleu}{HTML}{1F3A5F}
\definecolor{accent}{HTML}{C8892B}
\definecolor{fond}{HTML}{F4F6F9}

\setlength{\parindent}{0pt}
\setlength{\parskip}{0.45em}
\setlist[itemize]{label={\color{accent}\textbullet},leftmargin=1.4em,itemsep=0.2em}

% --- Titres -------------------------------------------------------------
\titleformat{\subsection}{\large\bfseries\color{bleu}}{}{0pt}{}
\titlespacing*{\subsection}{0pt}{0.9em}{0.25em}
\titleformat{\subsubsection}{\normalsize\bfseries\color{bleu!85}}{}{0pt}{}
\titlespacing*{\subsubsection}{0pt}{0.9em}{0.2em}

% En-tête d'une fiche ou d'une partie : petite étiquette + grand titre
\newcommand{\entete}[2]{%
  \clearpage
  {\color{accent}\small\bfseries\MakeUppercase{#1}}\par\vspace{0.1em}
  {\color{bleu}\LARGE\bfseries #2\par}
  \vspace{0.3em}{\color{accent}\rule{2.5cm}{2pt}}\par\vspace{0.6em}
  \addcontentsline{toc}{section}{#1 — #2}
  \markboth{#1}{#1}}

% --- En-tête et pied de page -------------------------------------------
\pagestyle{fancy}
\fancyhf{}
\fancyhead[L]{\small\color{bleu}Agôn Cup — R5A5, sujet B}
\fancyhead[R]{\small\color{gray}\leftmark}
\fancyfoot[C]{\small\color{gray}\thepage}
\renewcommand{\headrulewidth}{0.4pt}
\renewcommand{\headrule}{{\color{gray!40}\hrule height\headrulewidth}}

% --- Encadrés -----------------------------------------------------------
\newtcolorbox{prompt}{enhanced,breakable,colback=fond,colframe=fond,
  borderline west={3pt}{0pt}{bleu},arc=0pt,boxrule=0pt,
  left=10pt,right=8pt,top=6pt,bottom=6pt,fontupper=\small,
  before skip=0.5em,after skip=0.8em}
\newtcolorbox{extrait}{enhanced,colback=accent!8,colframe=accent!8,
  borderline west={3pt}{0pt}{accent},arc=0pt,boxrule=0pt,
  left=10pt,right=8pt,top=6pt,bottom=6pt,fontupper=\itshape,
  before skip=0.5em,after skip=0.8em}
\newtcolorbox{blocode}{enhanced,breakable,colback=fond,colframe=gray!30,boxrule=0.4pt,
  arc=2pt,left=6pt,right=6pt,top=4pt,bottom=4pt,
  before skip=0.5em,after skip=0.8em}
\fvset{fontsize=\footnotesize}

% Table des matières sobre
\makeatletter
\renewcommand{\l@section}[2]{\par\noindent{#1}\dotfill{#2}\par\vspace{0.35em}}
\makeatother

\begin{document}

% --- Page de titre ------------------------------------------------------
\thispagestyle{empty}
\begin{tikzpicture}[remember picture,overlay]
  \fill[bleu] (current page.north west) rectangle ([yshift=-7cm]current page.north east);
  \fill[accent] ([yshift=-7cm]current page.north west) rectangle ([yshift=-7.25cm]current page.north east);
\end{tikzpicture}
\vspace*{0.4cm}
{\color{white}
{\small\bfseries\MakeUppercase{R5A5 — Programmation avancée · Sujet B}}\par\vspace{0.8em}
{\fontsize{30}{36}\selectfont\bfseries Agôn Cup\par}\vspace{0.5em}
{\Large Dossier de décisions et rapport sur l'usage de l'IA\par}
}
\vspace{3.8cm}

{\large\color{bleu}\bfseries Denis Zapoi \quad·\quad Fabio Guerreiro-Marques\par}
\vspace{0.3em}
{\color{gray}BUT Informatique — IUT de Montreuil, Université Paris~8\par
Année 2026-2027 · Octobre 2026\par}

\vspace{2cm}
{\color{accent}\small\bfseries\MakeUppercase{Sommaire}}\par\vspace{0.6em}
\begingroup\color{bleu}
\makeatletter\@starttoc{toc}\makeatother
\endgroup


% =====================================================================
\entete{Fiche de décision n\textsuperscript{o}\,1}{Base de données}
\enlargethispage{4\baselineskip}

\subsection*{Le problème}

Pour commencer, il fallait réfléchir à la question SQL ou NoSQL, et la question a été vite répondue : au vu du projet, le relationnel semblait bien plus adapté. En effet, on a beaucoup de contraintes d'intégrité :
\begin{itemize}
  \item un joueur dans une seule équipe par tournoi ;
  \item un seul capitaine par équipe ;
  \item une seule demande d'adhésion en attente à la fois.
\end{itemize}

\subsection*{Les options}

Ensuite, il fallait trancher entre MySQL et PostgreSQL.

De plus, au départ, on avait envisagé d'utiliser Supabase (qui utilise PostgreSQL). On voulait choisir Supabase parce qu'on pensait, au début, que nous devions déployer le projet en ligne par la suite, et Supabase aurait hébergé la base de données. Cependant, notre professeur nous a dit qu'il allait simplement lancer le projet en local sur sa machine, le front comme le back, pour le regarder. Notre idée de départ, déployer le site avec une base hébergée en ligne, n'avait donc plus lieu d'être. Même si, au final, on n'a pas utilisé Supabase, c'était un argument de plus qui nous a fait pencher vers PostgreSQL.

\subsection*{Ce qu'on a choisi}

PostgreSQL, lancé avec Docker.

Un élément nous a fait aller vers PostgreSQL plutôt que MySQL : les index uniques partiels, qui servent par exemple à garantir l'unicité du capitaine. Sans cela, il aurait fallu vérifier ça uniquement au niveau du code.

\subsection*{Ce que ça coûte}

Au début, pas forcément de difficulté, puisqu'on codait directement sur les PC de l'IUT, donc tout allait bien. Pour simuler la base de données, on faisait des dumps. Ça a commencé à devenir pénible quand on a vu qu'on allait devoir implémenter une grosse base : il y avait les équipes, les tournois, les jeux, les rôles, les joueurs, les visiteurs, avec en plus des données sur les scores et l'avancée dans le tableau. Nos jeux de données devenaient vraiment énormes. L'idée de passer sur Docker vient de Fabio, et on s'est mis d'accord assez vite avec Denis. Avec le recul, c'était une bonne idée.

L'installation en elle-même nous a pris une dizaine de minutes. Le plus long a été de comprendre pourquoi ça ne marchait pas au départ : on a dû faire nos recherches pour vérifier que, sous Windows, il fallait bien installer WSL et Docker Desktop. On l'a fait sur nos deux ordinateurs.

Un autre coût est venu de Windows. On l'a vu tout de suite, parce que c'est Windows Sécurité qui nous affichait directement les blocages, et ça nous a posé pas mal de problèmes. Son contrôle d'application bloquait les exécutables générés dans l'environnement virtuel Python (comme \texttt{pip.exe}), et aussi la bibliothèque native embarquée dans psycopg, le pilote PostgreSQL habituel. Plutôt que de désactiver une protection de Windows, on a changé de pilote pour pg8000, qui est écrit entièrement en Python et n'a donc rien à bloquer. Il est plus lent que psycopg sur de gros volumes, ce qui n'a pas d'importance pour des tournois de quelques équipes, et le changement s'est limité au début de l'adresse de connexion à la base. On l'a noté dans le README pour ne pas bloquer quelqu'un d'autre.

% =====================================================================
\entete{Fiche de décision n\textsuperscript{o}\,2}{Mécanisme de temps réel}

\subsection*{Le problème}

Le sujet demande que les membres d'une équipe voient les messages arriver sans recharger la page, et que le suivi d'un tournoi se fasse en direct. Le chat est la partie du temps réel qui nous a le plus gênés. Ce n'est pas un chat global : c'est un chat par équipe. Dans un tournoi, il y a plusieurs équipes, et chaque équipe a son propre chat. Il fallait donc que seules les personnes présentes dans une équipe puissent communiquer entre elles, avec un cloisonnement entre les équipes.

\subsection*{Les options}

Le cours propose plusieurs mécanismes pour le temps réel. Voici ceux qu'on a regardés, avec les raisons pour lesquelles on les a gardés ou non :
\begin{itemize}
  \item \textbf{Le polling.} On n'a pas essayé : on a vu directement dans le cours que c'était une mauvaise idée, et on l'a donc évité tout de suite. Avoir autant de requêtes qui ne servent à rien, à demander toutes les $n$~secondes \og est-ce que tu as reçu un message ? \fg{}, ne semblait pas approprié au projet.
  \item \textbf{Le SSE} aurait été très approprié, et aurait totalement pu marcher. Pourquoi ne pas l'avoir utilisé ? Contrairement à Socket.IO, il n'y a pas de notion de salon (\emph{room}), donc il aurait fallu écrire nous-mêmes quel utilisateur peut écouter quel tournoi ou quelle équipe.
  \item \textbf{Le WebSocket} : même remarque sur la question des salons.
  \item \textbf{Socket.IO} : les salons sont natifs et la reconnexion est automatique.
\end{itemize}

\subsection*{Ce qu'on a choisi}

Socket.IO, avec la bibliothèque Flask-SocketIO côté back. Les salons ont permis de créer le cloisonnement dont on avait besoin pour le chat : un salon par équipe, et les personnes présentes dans cette équipe peuvent communiquer entre elles. Il propose directement la notion de salon, ce qui allège le travail qu'on doit réaliser.

\subsection*{Ce que ça coûte}

Un premier coût, tout bête, a été le test. Pour voir le temps réel, il faut deux comptes différents en même temps. En ouvrant deux pages du même navigateur, la session est partagée : se déconnecter sur une page déconnecte l'autre. On ne le savait pas, et ça nous a bloqués plusieurs jours. Au début, Fabio pensait que le problème venait de son code : qu'il avait mal implémenté quelque chose, ou mal compris une partie du cours. Il a refait certaines parties plusieurs fois, en se disant que c'était peut-être la déconnexion qui posait problème, mais ça ne venait pas de là. La solution était simplement d'utiliser la navigation privée pour avoir deux comptes en parallèle. C'était frustrant de rester bloqué aussi longtemps sur une erreur aussi bête, mais ça fait partie du métier de développeur : on se trompe, et on apprend de ses erreurs.

Un autre coût vient de la reconnexion. Socket.IO se reconnecte automatiquement quand le réseau coupe, mais pas quand c'est le serveur lui-même qui ferme la connexion (par exemple quand la session a expiré) : il considère alors la fermeture comme voulue. Le front a donc dû rouvrir la connexion lui-même dans ce cas.

Le coût le plus important a été la réflexion sur la sécurité. Dès qu'on touche à la sécurité, il faut penser à absolument tout, même à des choses qui ne semblent pas dangereuses au premier abord. Nos bases étaient limitées : on avait fait seulement quelques cours de sécurité pour comprendre le principe. On n'avait pas d'idée précise de la façon de faire, et on a donc dû beaucoup demander à l'IA, en lui demandant par exemple si une faille pouvait être laissée telle quelle. Elle répondait avec des exemples d'attaques et de conséquences qui nous ont fait corriger rapidement certains points.

Ce qui nous a le plus surpris, c'est une faille dans Werkzeug, la bibliothèque sur laquelle repose Flask (CVE-2026-102598). La semaine d'avant, nos versions de bibliothèques étaient bonnes. Entre-temps, une mise à jour de sécurité était sortie et on ne l'avait pas faite. Il a suffi de passer de la version 3.1.8 à la 3.1.9, et notre code n'utilisait pas la fonction concernée, mais on a été très surpris de découvrir une faille dans notre projet alors qu'on n'avait rien changé. Côté front, l'audit de Denis a aussi relevé deux alertes élevées sur Vite, corrigées de la même façon, par une mise à jour.

% =====================================================================
\entete{Fiche de décision n\textsuperscript{o}\,3}{Format et destination de la journalisation}
\enlargethispage{4\baselineskip}

\subsection*{Le problème}

On voulait retrouver dans les journaux les traces d'erreurs, le code de réponse de chaque requête, la date, et un message disant si l'opération avait réussi ou non. On voulait aussi y retrouver les refus d'accès et les tentatives d'infraction, par exemple quelqu'un qui essaie d'entrer là où il n'a pas le droit : en gros, tout ce qu'on attend de base de journaux. Chaque requête laisse la route, le code de réponse, la durée, l'utilisateur et un identifiant de requête.

\subsection*{Les options}

Le cours cite plusieurs destinations : la console, des fichiers avec rotation, une base de données, des plateformes centralisées et des services cloud.
\begin{itemize}
  \item \textbf{Une base de données.} On a écarté cette solution. La base de l'application a une structure faite pour évoluer, et elle contient des données qu'on veut protéger. Créer une autre base dédiée uniquement à la journalisation aurait été long.
  \item \textbf{Des fichiers.} Il faut les faire tourner, les purger et surveiller le disque, et ils disparaissent avec le conteneur.
  \item \textbf{Une plateforme centralisée ou un service cloud.} C'est lourd à installer et à maintenir, et prévu pour de nombreux services, pas pour une seule application.
\end{itemize}

\subsection*{Ce qu'on a choisi}

La sortie standard, avec une ligne JSON (JSON Lines) par événement. Pour la destination, le cours préconise la console quand l'application tourne dans un conteneur, ce qui est notre cas :
\begin{extrait}
\og En conteneur, la console redevient la bonne destination. [\ldots] L'application se contente de produire du JSON Lines et ignore où il finit. \fg{}
\end{extrait}
Le niveau suit le résultat : un succès est en INFO, un refus en WARNING, une erreur du serveur en ERROR.

\subsection*{Ce qu'on journalise, et ce qu'on ne journalise pas}

On ne journalise ni mot de passe, ni jeton, ni e-mail, ni contenu des requêtes, parce que ce sont des informations personnelles. Même si on a accès au serveur, on n'est pas censés les garder : quelqu'un qui aurait accès à une console verrait aussi les journaux, donc les mots de passe. Au début du développement, on avait mis des mots de passe et des informations de ce genre dans les journaux pour comprendre les erreurs. À la toute fin, on a vérifié qu'on avait tout retiré, et on n'a retrouvé aucun mot de passe oublié.

\subsection*{Ce que ça coûte}

Le plus gênant, c'est que les journaux disparaissent quand le conteneur est supprimé. Pour être sûrs que le projet marche pour quelqu'un d'autre et que le README est à jour, on le clonait et on repartait de zéro, comme pour une première installation. On avait alors beaucoup de copies de journaux à garder pour comparer les anciens avec les nouveaux.

L'autre coût est la lecture. Dans d'autres outils, comme Log4j en Java, on peut mettre en page les journaux avec des couleurs. Avec notre pile, on ne savait pas comment faire, et nos premières lignes de JSON étaient difficiles à lire, parce qu'on n'avait pas bien configuré leur mise en page. Le projet a un mode texte (\texttt{LOG\_FORMAT=texte}) plus lisible en développement.

% =====================================================================
\entete{Rapport}{L'usage de l'IA dans le projet}

\subsection*{1. Ce que nous avons utilisé}

On a surtout utilisé Claude, parce que l'un de nous (Fabio) a un abonnement Claude. L'outil principal a été Claude Code, utilisé dans des sessions en ligne reliées à GitHub, qui ouvrait les pull requests. On n'a pas configuré de skills, de sous-agents ni de revue automatisée.

Pour la base de données, on avait proposé plusieurs choix et on a demandé à Claude lequel était le meilleur par rapport à notre projet, en lui donnant le contexte, et s'il existait d'autres technologies. Pour ses audits de sécurité, l'un de nous (Denis) a aussi utilisé Claude.

Tout est visible dans les dépôts : les commits faits avec Claude Code y sont signés. Sur la branche principale, Claude Code a signé 20 commits sur 41 dans le back, et 24 sur 40 dans le front.

\textbf{Les deux fichiers \texttt{CLAUDE.md}.} On a des cours de programmation multimédia, qui sont des cours de programmation avec l'IA, et aujourd'hui, si on veut être à jour et rapide, on est obligé d'utiliser l'IA. Quand une personne ouvre notre dépôt avec Claude Code, Claude lit automatiquement \texttt{CLAUDE.md} et \texttt{API.md} : il a déjà l'architecture, une idée du fonctionnement du projet et les contrats entre le back et le front. Cela permet aux prochaines personnes qui voudront utiliser l'IA sur notre projet d'avoir une IA qui le comprend déjà très bien. Ces fichiers contiennent aussi les règles de travail qu'on a données à Claude : parler français, faire une branche et une pull request par changement, ne jamais fusionner soi-même, ne jamais mettre de secret dans le dépôt, et ne pas rédiger à notre place les fiches de décision ni ce rapport.

\subsection*{2. Nos prompts}

Voici trois de nos prompts, recopiés tels quels, avec nos fautes de dictée.

\subsubsection*{Prompt 1 — partir d'un plan, en validant chaque étape (15 octobre)}

\begin{prompt}
Ok, bon, j'ai pas encore accepté les PR, mais on fera ça plus tard. J'aimerais que tu analyses du coup mon projet, donc mon back et mon front. Et j'aimerais, parce qu'en fait, tu vois, pour le moment, je n'ai que créé la page de connexion ainsi que euh, la page principale. Euh, maintenant, ce que j'aimerais bien faire, c'est voilà, tu vois, commencer à m'attarder à bah, création de tournoi, euh, gestion du tournoi, euh, système de chat et compagnie. Genre vraiment, je vais pas dire finir le projet, tu vois, mais vraiment avancer, avancer, avancer. Est-ce que tu pourrais me fournir euh, des documents, un document, un, soit un point MD ou n'importe, hein, je ne sais pas. Des classes que je dois rajouter, des trucs que je dois modifier, compagnie. Donc, je te demanderai pas toi de me faire les les modifications. Je le ferai moi-même. Mais voilà, dis-moi euh, quelles sont les classes que je dois changer, les éléments. Enfin, donne-moi le code en entier de chaque classe. Voilà, euh, de celles qui sont modifiées ou supprimées ou ajoutées ou n'importe. Mais comme ça, ce sera moi qui fera le commit et compagnie, tu vois. Donc, est-ce que tu peux lire le sujet, regarder les attentes, et à partir de ce moment-là, et à partir de ce que j'ai déjà fait et compagnie, voilà, commencer déjà à mettre en place le système de tournoi, le système de chat, de compte, euh, et compagnie, quoi. J'ai envie de te dire un peu avoir déjà fini le projet. On fera ça petit à petit. Si tu peux pas faire ça en une seule demande, ce n'est pas grave. On va faire petit à petit, comme ça on peut valider chaque étape.
\end{prompt}

\textbf{Ce que ça donnait au départ.} À ce moment, on n'avait fait que la page de connexion et la page principale. On voulait que Claude nous dise quelles classes ajouter ou modifier, avec leur code entier, pour faire nous-mêmes les commits.

\textbf{Ce qu'on a changé.} Le prompt demandait d'avancer par petites étapes, validées une à une. Plus tard, on est passés de \og je fais les modifications moi-même \fg{} à \og corrige dans des pull requests \fg{} (prompt~3).

\textbf{Ce que ça a donné.} Claude a livré la première étape, les équipes et les demandes d'adhésion côté back : le code complet de 22~fichiers, un plan général en cinq étapes, et 75~nouveaux tests (156 en tout). Il a aussi signalé un défaut : \texttt{docker compose up} aurait échoué sur une base neuve, parce que les données de démonstration étaient jouées avant le schéma.

\subsubsection*{Prompt 2 — cadrer la question pour obtenir une réponse courte (15 octobre)}

\begin{prompt}
dis, j'ai juste une question comme ça. Tu sais, on a vu le temps réel et compagnie et tout. Je me disais, est-ce que dans le système de chat, est-ce qu'on ferait un truc qui permettrait d'envoyer des photos ou des vidéos ou compagnie, ou est-ce que ce serait trop compliqué, il faudrait gérer ça côté serveur ou des trucs comme ça ? Parce qu'après, en fait, le rendu, qu'on doit donner, c'est un truc que notre prof, il va déployer lui-même chez lui et tout, tu vois, genre pour tester. Donc je me dis, est-ce qu'on va vraiment, on a besoin de faire ça ? Est-ce qu'on peut faire ça ou pas ? Je veux juste une réponse. Je veux un oui ou un non avec deux, trois petites explications, mais voilà, commence pas à partir dans les détails ou commence à vouloir implémenter ça.
\end{prompt}

\textbf{Ce que ça donnait au départ.} Une idée en passant : permettre d'envoyer des photos ou des vidéos dans le chat.

\textbf{Ce qu'on a changé.} On a précisé qu'on voulait seulement une réponse oui ou non, avec deux ou trois explications, sans partir dans les détails ni commencer à implémenter, pour garder le contrôle de ce qui était fait.

\textbf{Ce que ça a donné.} Claude a répondu non : le sujet demande un chat d'équipe, pas l'envoi de fichiers ; stocker des fichiers demande de vérifier leur type et leur taille et d'appliquer à chacun la même règle d'accès que pour les messages ; et plus il y a de choses à configurer, plus il y a de risques que le projet ne démarre pas chez le professeur. On n'a pas implémenté l'envoi de fichiers.

\subsubsection*{Prompt 3 — déléguer des corrections dans des pull requests (27 octobre)}

\begin{prompt}
Ok, bah écoute, s'il te plaît, est-ce que tu peux lister les dernières, enfin, corriger les dernières choses qui manquent à travers une PR que tu feras dans le back et dans le front. Ensuite, après, tu me parles d'un écran, mon prochain match, puis un Docker Compose pour toute la pile. Ajoute ça aussi, parce que c'est très important, par la même occasion. Tu m'as dit, du coup, que la documentation et tout ça a été faite. Bah écoute, je te crois à ce niveau-là. Bon, du coup, je vais éviter de t'envoyer chaque cours de chaque page de Thomas Louvet parce que c'est un peu énorme. Enfin, la page de mon prof, c'est un peu beaucoup. Mais du coup, je pense que tu peux juste te baser au niveau des chapitres et ça devrait suffire. La manière dont on va faire ça, c'est voilà. Fais des PR pour faire des modifications par rapport à ce qui manque. Ensuite, après avoir fait tout ça, on va changer 2-3 trucs, toi et moi, au niveau du site, genre le texte et compagnie et tout. Et ensuite, après avoir fait ça, on va s'assurer de faire des jolis readme pour un bon déploiement, pour un déploiement facile en local sur l'ordinateur de mon prof. Ensuite, après avoir fait ça, on va s'assurer\ldots{} encore une fois de la documentation hein, de façon à ce que ça soit bon et ensuite après on va faire tout le travail de journalisation donc voilà.
\end{prompt}

\textbf{Ce que ça donnait au départ.} Une demande large : corriger ce qui manquait, ajouter l'écran \og mon prochain match \fg{} et un docker-compose pour toute la pile, puis passer aux README et à la journalisation.

\textbf{Ce qu'on a changé.} Contrairement au prompt~1, on ne faisait plus les modifications nous-mêmes : elles passaient par des pull requests faites par Claude, que l'on relisait et fusionnait ensuite. On a aussi décidé de ne pas envoyer chaque page du cours, mais de laisser Claude s'appuyer sur les chapitres.

\textbf{Ce que ça a donné.} Des pull requests dans les deux dépôts, dont un \texttt{docker-compose.complet.yml} pour démarrer la base, l'API et l'interface ensemble, et l'écran \og mon prochain match \fg{}.

\subsection*{3. Ce que nous avons corrigé}

Le code produit par l'IA n'était pas toujours juste, et on a dû le reprendre.

\subsubsection*{Une course dans l'abonnement à une équipe (Denis)}

Lors de l'audit de sécurité de Denis, on a trouvé que la vérification des droits avant d'entrer dans le salon d'une équipe laissait une petite fenêtre : si l'exclusion d'un joueur était validée entre la vérification et l'entrée dans le salon, le joueur exclu entrait quand même et recevait les messages de l'équipe. Un test a été écrit pour reproduire le problème : il échouait avant le correctif. Voici le code livré à la place, qui revérifie les droits après l'entrée :

\begin{blocode}
\begin{Verbatim}
# Revérifié APRÈS l'entrée : une exclusion validée entre la première
# vérification et join_room a retiré le joueur d'un salon où il n'était
# pas encore, puis il y serait entré.
if not espace_service.autoriser_abonnement(id_equipe, identite["id"], identite["ver"]):
    leave_room(salon_equipe(id_equipe))
    metriques.abonnement_refuse()
    return {"ok": False}
\end{Verbatim}
\end{blocode}

\subsubsection*{Le badge \og non~lu \fg{} qui disparaissait au rechargement}

En vérifiant le chat à deux comptes, Claude a rejoué le scénario dans un vrai navigateur et a vu qu'après un rechargement de la page, un message arrivé pendant l'absence n'était plus signalé \og non~lu \fg{}. Au rechargement, la session est restaurée après le premier affichage : la messagerie lisait donc ce qui avait déjà été vu avec un identifiant encore vide, et la première lecture des messages écrasait la mémoire enregistrée dans le navigateur. On a écrit un test qui échoue sans le correctif, puis corrigé ainsi :

\begin{blocode}
\begin{Verbatim}
// Avant : la mémoire est lue une seule fois, au premier affichage.
const [lus, setLus] = useState(() => lireLus(moiId));

// Après : elle est relue quand l'identifiant de l'utilisateur change.
const [memoire, setMemoire] = useState(() => ({ pour: moiId, lus: lireLus(moiId) }));
if (memoire.pour !== moiId) setMemoire({ pour: moiId, lus: lireLus(moiId) });
\end{Verbatim}
\end{blocode}

\subsubsection*{D'autres corrections}

\begin{itemize}
  \item \textbf{Le cookie de session} durait un an (valeur par défaut de la bibliothèque), alors que le jeton expire au bout de 12~heures. Ce défaut a été relevé quand on a demandé à Claude de comparer le projet au sujet, puis corrigé en fixant la durée du cookie.
  \item \textbf{Les autres points de l'audit de Denis} : un joueur pouvait modifier le nom de son équipe après la clôture des inscriptions, une écriture venant d'un autre site pouvait être acceptée, et la limite d'inscription était trop large (5 par minute au lieu de 10 par heure).
  \item \textbf{La documentation} comportait des erreurs, trouvées lors d'un audit final : un numéro de version du contrat périmé, et un tableau qui disait le contraire du code sur l'accès des administrateurs aux salons d'équipe.
\end{itemize}

\subsection*{4. Bilan}

\textbf{Ce que ça nous a apporté.} Pour le front, on avait les cours de M.~Thomas Louvet, à qui on rend ce projet, donc on avait des bases. Pour le back, c'était différent : Fabio n'avait eu qu'un cours sur Flask et sur les API, et les alternants n'avaient pas eu grand-chose non plus. Denis avait beaucoup plus d'expérience, et Fabio a beaucoup appris de lui sur la façon de coder avec Flask et le fonctionnement d'une API. L'IA l'a aussi aidé en parallèle, en lui expliquant ce qu'il ne comprenait pas, souvent tard le soir.

\textbf{Ce que ça nous a coûté.} Beaucoup d'investissement. On a dû prendre des notes en continu sur ce qu'on faisait, même les choses les plus minimes, pour pouvoir écrire le document que vous avez sous les yeux. Ça a pris du temps, et c'était assez pénible. En plus, les perturbations des dernières semaines nous ont ralentis : avec les manifestations étudiantes autour du CROUS et de l'Université Paris~8, et les blocus, il nous est arrivé de venir à l'IUT en se disant qu'on allait pouvoir travailler, et de ne pas pouvoir étudier. Ce n'étaient pas du tout des conditions favorables à l'apprentissage. Au final, le projet a représenté beaucoup de travail à la maison, sur des heures prévues pour d'autres projets ou d'autres choses.

\textbf{Ce qu'on ferait autrement.} Honnêtement, on ne voit pas vraiment ce qu'on aurait fait autrement. On aurait juste voulu que le projet se déroule dans d'autres conditions.

\end{document}
